\documentclass{article}
\usepackage[T1]{fontenc}
\usepackage{lmodern}
\usepackage{graphicx}
\usepackage{amsmath,amssymb}
\usepackage[a4paper,margin=2cm]{geometry}
\usepackage[absolute,overlay]{textpos}
\setlength{\TPHorizModule}{1mm}
\setlength{\TPVertModule}{1mm}
\usepackage[indonesian]{babel}
\usepackage{hyperref}
\setlength{\emergencystretch}{2em}
% unit: Halaman 49

\begin{document}

% === Halaman 49 (python/native) ===

• Berikutnya, dekripsi C dengan semua semua 

kemungkinan nilai K2 (yaitu sebanyak 256 

kemungkinan kunci).

• Bandingkan semua hasil dekripsi ini dengan 

elemen di dalam tabel tadi. Jika ada yang 

sama, maka dua buah kunci, K1 dan K2, telah 

ditemukan. 

• Tes kedua kunci ini dengan pasangan plainteks-

cipherteks lain yang diketahui. Jika kedua kunci 

tersebut menghasilkan cipherteks atau 

plainteks yang benar, maka K1 dan K2 tersebut 

merupakan kunci yang benar  

49

K1             X                            K2           X

X = EK1(P)

X = DK2(C)

000…00       1010..11

000…01       1101..00

…

…

…

101..10         010..01

…

…

… 

…

...

111..11        011.. 10

000…00       0110..10

000…01       1000..10

…

…

…

101..10         010..01

…

…

110..10…     010..01 

…

...

111..11        011.. 10

\end{document}
